%% Information and Software Technology research article draft.
%% Prepared against the IST Guide for Authors accessed 2026-07-15.
\documentclass[preprint,12pt,a4paper]{elsarticle}

\usepackage{amssymb}
\usepackage{amsmath}
\usepackage{booktabs}
\usepackage{array}
\usepackage{enumitem}
\usepackage{float}
\usepackage{hyperref}
\usepackage{url}

\hypersetup{
  colorlinks=true,
  linkcolor=blue!55!black,
  citecolor=blue!55!black,
  urlcolor=blue!55!black
}
\urlstyle{same}
\emergencystretch=3em
\setlist[itemize]{leftmargin=1.5em}
\setlist[enumerate]{leftmargin=1.5em}
\journal{Information and Software Technology}

\begin{document}

\begin{frontmatter}

\title{Evidence-Closed Knowledge Maintenance for LLM Agents: Design and Empirical Evaluation of Agent Wiki}

\author[independent]{Zijun Chu\corref{cor1}}
\ead{chunimav5@gmail.com}
\ead[url]{https://github.com/NimaChu/agent-wiki}
\cortext[cor1]{Corresponding author.}
\address[independent]{Independent Researcher, Room 1303, No. 19, Lane 899,
Xinfeng Middle Road, Qingpu District, Shanghai, China}

\begin{abstract}
\textbf{Context:} Large language model (LLM) agents need durable external
knowledge, but file-based corpora provide little assurance that source evidence
remains reachable after agents synthesize or reorganize it.
\textbf{Objective:} This study investigates whether an executable
evidence-closure invariant can make local Markdown knowledge maintenance
structurally auditable at documentation scale.
\textbf{Methods:} We designed and implemented Agent Wiki, a local-first system
that separates raw evidence from synthesized wiki pages and checks four closure
conditions for processed sources. We evaluated release v0.2.1 using a public
FlexSim documentation-derived graph containing 1,092 raw records, 88 wiki
records, and 5,163 edges. We then injected 120 faults from four maintenance
fault classes into a 1,000-record fixture and compared Agent Wiki with an
existence-only link baseline. Finally, we measured end-to-end scan and closure
checking over synthetic vaults of 100--2,000 raw records, with ten repetitions
per size.
\textbf{Results:} All 1,092 public raw records retained a source URL, a SHA-256
content hash, a resolved related wiki target, and a reciprocal wiki-to-evidence
backlink. Agent Wiki detected all 120 injected faults with no false positives;
the existence-only baseline detected 30. Median end-to-end checking time ranged
from 56.8 ms for 100 raw records to 1,302.3 ms for 2,000.
\textbf{Conclusion:} Executable evidence closure detects lifecycle failures
beyond broken links while retaining low local checking cost at the studied
scale. The results establish structural integrity, not semantic correctness or
superior answer quality relative to retrieval-augmented generation.
\end{abstract}

\begin{keyword}
LLM agents \sep knowledge maintenance \sep traceability \sep provenance
\sep local-first software \sep Markdown \sep empirical software engineering
\end{keyword}

\end{frontmatter}

\section{Introduction}

Large language model (LLM) agents increasingly perform multi-step software and
knowledge-work tasks. They inspect repositories, read technical documentation,
call external tools, modify files, and resume work across sessions. Systems and
benchmarks such as ReAct, SWE-agent, and SWE-bench demonstrate the value of
combining language models with actions and software environments
\cite{yao2023react,jimenez2024swebench,yang2024sweagent}. These workflows also
expose a less visible engineering problem: the knowledge accumulated by an
agent must remain inspectable and maintainable after the immediate conversation
ends.

External memory is commonly implemented through chat history, application
memory, document retrieval, vector indexes, or human-authored notes. Retrieval-
augmented generation (RAG) and graph-based retrieval improve access to external
information \cite{lewis2020rag,edge2024graphrag}. Memory-oriented systems such
as MemGPT manage information across context tiers \cite{packer2023memgpt}.
These approaches address what information should be presented to a model, but
they do not by themselves define when a captured source has been responsibly
integrated into a durable knowledge base. A source may be marked processed even
though its synthesis page is missing, its target link is broken, the synthesis
does not point back to evidence, or an explicit follow-up remains open.

This distinction is a software-quality concern. For a corpus maintained by
humans and agents, structural consistency must be checkable independently of
the model that created the notes. Plain files provide ownership and
inspectability, but file ownership alone does not provide lifecycle invariants.
A conventional broken-link checker can identify a reference to a missing file;
it cannot determine that a processed source declares no durable target, that a
wiki page fails to preserve a backlink to evidence, or that processing was
closed while a follow-up flag remained active.

We address this problem with \emph{evidence closure}: an executable invariant
over a two-layer Markdown knowledge base. Raw notes preserve evidence and
capture metadata; wiki pages contain durable synthesis. A raw source may be
treated as processed only when it declares at least one related wiki target,
those targets resolve, at least one related wiki page links back to the raw
source, and no explicit follow-up remains. Agent Wiki implements this invariant
as part of a local command-line workflow and exposes the same derived graph to
search, maintenance queues, and an optional read-only dashboard.

The article makes four contributions:

\begin{enumerate}
  \item an operational evidence-closure model for agent-maintained knowledge,
  expressed as conditions that can be checked over ordinary Markdown files;
  \item an open-source implementation that separates evidence, synthesis, and
  derived visualization while requiring no hosted database or embedding
  service;
  \item a reproducible empirical evaluation combining a public documentation-
  derived graph, seeded maintenance faults, and an existence-only baseline;
  and
  \item a runtime study over controlled vaults containing up to 2,000 raw
  records, together with scripts and machine-readable results.
\end{enumerate}

The evaluation is intentionally scoped to structural integrity. It does not
claim that evidence closure proves factual correctness, improves user
productivity, or outperforms RAG on answer quality. Instead, it asks whether a
specific, reviewable maintenance property can be enforced and measured.

\section{Background and Related Work}

\subsection{LLM Agent Memory and Retrieval}

RAG combines parametric generation with non-parametric document retrieval and
has become a standard architecture for knowledge-intensive language tasks
\cite{lewis2020rag}. GraphRAG constructs entity graphs and community summaries
to support global questions over large corpora \cite{edge2024graphrag}.
MemGPT treats context management as a tiered memory problem, moving information
between constrained and external storage \cite{packer2023memgpt}. These systems
show that external knowledge is central to useful long-running agents.

Agent Wiki operates at a different layer. It does not prescribe an answer-time
retriever, chunking strategy, embedding model, or prompt. It maintains the
inspectable corpus from which a retriever may later be built. This separation
allows vector, lexical, graph-guided, or manual retrieval to share the same
source and synthesis artifacts. It also allows integrity checks to run without
an LLM or network service.

\subsection{Local-First and Personal Knowledge Systems}

Local-first software emphasizes ownership, offline availability, and user
control over data \cite{kleppmann2019localfirst}. Markdown knowledge tools such
as Obsidian and Logseq demonstrate the practical value of file-native notes,
backlinks, and graph views \cite{obsidian,logseq}. Evergreen-note practices
similarly encourage small, durable notes that evolve through linking and
revision \cite{matuschakEvergreen}.

These tools make knowledge inspectable, but their generality leaves lifecycle
policy to the user. Agent Wiki adds an agent-operating protocol: repository
rules define which files are raw evidence, which files are synthesis, which
metadata fields describe state, and which checks must pass before a source is
considered processed. The graph remains derived from the files rather than
becoming a second source of truth.

\subsection{Traceability and Provenance}

Traceability research has long treated links between artifacts as essential
for understanding origins, dependencies, and change \cite{gotel1994trace}.
The W3C PROV family provides a general model for representing entities,
activities, agents, and provenance relations \cite{w3cprov}. Knowledge graphs
provide a broader framework for representing entities and relations in a
machine-processable graph \cite{hogan2021knowledgegraphs}.

Evidence closure is narrower than general provenance modeling. It does not
attempt to encode every transformation or claim. It defines a minimum
bidirectional traceability gate between a captured source and at least one
durable synthesis artifact. The small scope makes the invariant easy to run in
local development and continuous integration. It also exposes a concrete
failure model that can be evaluated through mutation.

\subsection{Research Gap}

Prior work provides retrieval, memory management, local ownership, or general
provenance models. The missing practical element is a lightweight, executable
definition of \emph{processed} for an agent-maintained file corpus. Table
\ref{tab:positioning} summarizes this positioning. The categories are
architectural tendencies rather than claims about every implementation.

\begin{table}[H]
\centering
\small
\begin{tabular}{>{\raggedright\arraybackslash}p{0.21\linewidth}
  >{\centering\arraybackslash}p{0.13\linewidth}
  >{\centering\arraybackslash}p{0.13\linewidth}
  >{\centering\arraybackslash}p{0.20\linewidth}
  >{\centering\arraybackslash}p{0.13\linewidth}}
\toprule
Approach & Local files & Retrieval & Raw/synthesis split & Closure gate \\
\midrule
Chat transcript & Sometimes & Context search & No & No \\
Vector/RAG pipeline & Optional & Yes & Optional & No \\
Markdown note system & Yes & Optional & User-defined & User-defined \\
Agent Wiki & Yes & Optional & Yes & Yes \\
\bottomrule
\end{tabular}
\caption{Architectural positioning of Agent Wiki.}
\label{tab:positioning}
\end{table}

\section{Evidence-Closure Model}

\subsection{Two-Layer Knowledge Representation}

Let $R$ be the set of raw source notes and $W$ the set of synthesized wiki
pages. A raw note records capture metadata, including source URL, capture time,
content hash, source quality, status, and intended wiki targets. Its body may
contain captured text, extracted claims, image inventories, and processing
notes. A wiki page contains one durable knowledge unit, such as a concept,
method, API, workflow, entity, or comparison, and cites one or more raw notes.

Separating $R$ and $W$ serves two purposes. First, synthesis can change without
rewriting captured evidence. Second, a reader can distinguish what a source
contained from what an agent or human concluded from it. A single source may
support multiple wiki pages, and a wiki page may synthesize multiple sources.

\subsection{Closure Invariant}

For a raw source $r \in R$, let $\mathrm{related}(r)$ be its declared wiki
targets, $\mathrm{resolve}(t)$ resolve a target to a page or failure $\bot$,
and $\mathrm{links}(w)$ be the links emitted by wiki page $w$. The source is
evidence-closed only if:

\begin{equation}
\begin{split}
\mathrm{closed}(r) \equiv {} & |\mathrm{related}(r)| > 0 \\
& \land\ \forall t \in \mathrm{related}(r),\ \mathrm{resolve}(t) \neq \bot \\
& \land\ \exists t \in \mathrm{related}(r):
r \in \mathrm{links}(\mathrm{resolve}(t)) \\
& \land\ \neg\mathrm{needsFollowup}(r).
\end{split}
\label{eq:closure}
\end{equation}

The lifecycle rule is then:

\begin{equation}
\mathrm{status}(r)=\texttt{processed}\ \Rightarrow\ \mathrm{closed}(r).
\label{eq:gate}
\end{equation}

Violations are reported with four reason codes: \texttt{missing-related},
\texttt{unresolved-related}, \texttt{missing-wiki-backlink}, and
\texttt{explicit-followup}. These codes form the fault model used in the
evaluation.

\subsection{Scope of the Invariant}

Closure is necessary but not sufficient for semantic quality. A backlink can
exist while a synthesis is inaccurate. A content hash records a captured
state but does not prove the trustworthiness of the source. A resolved target
may be too broad or too narrow. Agent Wiki therefore reports closure as a
mechanical property and retains source-quality metadata for human or agent
judgment. Claim-level entailment, answer faithfulness, and source authority are
outside Equations \ref{eq:closure}--\ref{eq:gate}.

\section{System Design and Implementation}

\subsection{Architecture}

Agent Wiki is a self-contained repository implemented with Node.js scripts and
an optional Vite/React dashboard. Figure \ref{fig:architecture} shows the data
flow. Markdown remains the source of truth; graph JSON and the dashboard are
rebuildable views.

\begin{figure}[H]
\centering
\fbox{\begin{minipage}{0.90\linewidth}
\centering
Source capture and metadata\\
$\downarrow$\\
\textbf{Raw evidence layer} (captures, hashes, snapshots, image inventories)\\
$\updownarrow$ evidence links and closure checks\\
\textbf{Wiki synthesis layer} (concepts, methods, APIs, workflows)\\
$\downarrow$ derived scan\\
Search, lint, maintenance queues, graph JSON, and read-only dashboard
\end{minipage}}
\caption{Agent Wiki keeps evidence and synthesis in Markdown and derives
operational views from the same scanner.}
\label{fig:architecture}
\end{figure}

The main directories are \path|raw/| for evidence, \path|wiki/| for synthesis,
\path|templates/| for reusable schemas, \path|src/| for command-line logic,
and \path|tools/wiki-dashboard/| for visualization. Public software and
private knowledge can be separated through version-control policy.

\subsection{Shared Scanner and Resolver}

The scanner walks Markdown files, parses a deliberately constrained
frontmatter subset, extracts body and metadata wiki links, and resolves targets
by identifier, title, basename, or alias. It then emits nodes, edges, unresolved
links, typed relations, incoming and outgoing degree counts, and excerpts. The
same resolver supports status, lint, search, link repair, and dashboard graph
generation. Reusing one resolver prevents each interface from applying a
different definition of a valid link.

\subsection{Operational Workflows}

The command-line interface exposes capture, search, status, lint, garden review,
link repair, image indexing, and graph generation through root npm scripts.
The dashboard is read-only: write operations remain explicit Markdown edits or
command-line actions. This choice makes changes visible to ordinary diff and
review tools.

Image-rich sources receive a parallel evidence inventory. The image workflow
extracts Markdown and HTML image references, resolves source-relative URLs,
mirrors permitted images locally, and writes a machine-readable index. Images
are not part of the closure invariant evaluated here, but their metadata
preserves another evidence channel that text-only capture would lose.

\subsection{Implementation of Closure Checks}

After link resolution, the checker selects raw nodes whose status is
\texttt{processed}. It verifies the four conditions in Equation
\ref{eq:closure}. The check is deterministic and requires no model inference.
Reports are emitted as JSON so failures can be inspected by humans, agents, or
continuous-integration jobs. The current implementation favors portability and
reviewability over a complete YAML or Markdown dialect.

\section{Research Method}

\subsection{Research Questions}

The evaluation addresses three research questions:

\begin{description}
  \item[RQ1:] To what extent does a documentation-scale public case study
  satisfy the observable evidence-closure preconditions?
  \item[RQ2:] How effectively do evidence-closure checks detect seeded
  maintenance faults compared with an existence-only link baseline?
  \item[RQ3:] What is the end-to-end runtime cost of scanning and closure
  checking as the number of Markdown records increases?
\end{description}

\subsection{Artifact and Reproduction Package}

The study evaluates Agent Wiki release \texttt{v0.2.1}. The software is
licensed under the MIT License. Release \texttt{v0.2.1} and the public dataset
are archived with the current DOI. The experiment script, JSON results, and CSV
tables accompany this manuscript and are intended for the next archival release
before submission. The complete experiment is invoked from the repository root
with:

\begin{verbatim}
npm run paper:ist-evaluate
\end{verbatim}

The script reads the public dataset without modifying it, constructs controlled
fixtures in an operating-system temporary directory, runs the measurements,
writes machine-readable results beside the manuscript, and deletes temporary
fixtures.

\subsection{RQ1: Public Case-Study Dataset}

The public case study is derived from the English Autodesk FlexSim 2026 online
help captured for an authorized local knowledge base. To avoid redistributing
third-party documentation, the public dataset contains only whitelisted source
metadata, structural features, hashes, and graph edges. It excludes captured
text, webpage snapshots, screenshots, images, binaries, and local absolute
paths.

The dataset contains \texttt{raw-sources.jsonl}, \texttt{wiki-pages.jsonl},
\texttt{edges.jsonl}, a schema, and aggregate metrics. For each raw record, the
evaluation tests whether a source URL and 64-character SHA-256 hash are present,
whether at least one declared wiki target resolved, and whether at least one
resolved target has a reverse edge to the raw record. The last measure is the
public-data counterpart of reciprocal evidence closure.

\subsection{RQ2: Fault-Injection Protocol}

We generated a clean fixture with 1,000 processed raw notes and 40 synthesized
wiki pages. Each raw note mapped to one wiki page, and each wiki page linked back
to its evidence. We then injected 30 non-overlapping instances of each fault
class in Table \ref{tab:faultmodel}, for 120 total faults. Fault locations were
fixed by the script, making the run deterministic.

\begin{table}[H]
\centering
\small
\begin{tabular}{>{\raggedright\arraybackslash}p{0.27\linewidth}p{0.62\linewidth}}
\toprule
\normalfont Fault & Mutation \\
\midrule
\texttt{missing-related} & Remove the raw note's declared wiki target. \\
\texttt{unresolved-related} & Replace the target with a nonexistent wiki page. \\
\texttt{missing-wiki-}\allowbreak\texttt{backlink} & Remove the raw-source link from the related wiki page. \\
\texttt{explicit-followup} & Set \texttt{needs\_followup} to true on a processed source. \\
\bottomrule
\end{tabular}
\caption{Seeded maintenance fault model.}
\label{tab:faultmodel}
\end{table}

The comparison baseline represents a conventional existence-only Markdown link
checker: a fault is detected when an explicit link target does not resolve. It
does not interpret source status, require a target declaration, inspect reverse
evidence links, or enforce follow-up state. This baseline is intentionally
narrow; it isolates the additional behavior contributed by evidence closure
and is not a comparison with RAG or a database-backed knowledge system.

We report true positives, false positives, false negatives, precision, and
recall against the seeded fault labels. Because the experiment is deterministic
and exhaustively labels the constructed fixture, we use descriptive measures
rather than inferential tests.

\subsection{RQ3: Runtime Protocol}

We generated clean fixtures containing 100, 500, 1,000, and 2,000 raw records.
One wiki page was created for each group of at most 25 raw records, giving 4,
20, 40, and 80 wiki records. Every raw note was processed and evidence-closed.
The measured operation includes filesystem traversal, UTF-8 reads, frontmatter
and link parsing, link resolution, graph statistics, and processed-source
closure checking.

For each size, two warm-up runs preceded ten measured repetitions. We report
median, 95th percentile, minimum, and maximum wall-clock time. Measurements ran
on Windows 11 with Node.js 20.20.0, an Intel Core i9-14900K, 32 logical CPUs,
and 127.7 GiB of memory. The implementation is single-process for this workflow;
the hardware description supports replication but is not used to claim general
performance across machines.

\section{Results}

\subsection{RQ1: Observable Closure at Documentation Scale}

Table \ref{tab:dataset} reports the public case-study graph. All 1,092 raw
records contained a source URL and SHA-256 content hash. Every raw record had at
least one resolved related wiki target, and every raw record had a reciprocal
edge from at least one related wiki page. Thus, all observable closure
preconditions represented in the public dataset had 100\% coverage.

\begin{table}[H]
\centering
\small
\begin{tabular}{lr}
\toprule
Measure & Value \\
\midrule
Raw source records & 1,092 \\
Synthesized wiki records & 88 \\
Graph nodes & 1,180 \\
Graph edges & 5,163 \\
Raw-to-wiki edges & 2,273 \\
Wiki-to-raw edges & 2,310 \\
Wiki-to-wiki edges & 580 \\
Raw records with source URL & 1,092 (100\%) \\
Raw records with SHA-256 hash & 1,092 (100\%) \\
Raw records with resolved related target & 1,092 (100\%) \\
Raw records with reciprocal backlink & 1,092 (100\%) \\
\bottomrule
\end{tabular}
\caption{Public FlexSim case-study dataset, release v0.2.1.}
\label{tab:dataset}
\end{table}

This result answers RQ1 for the released dataset. It demonstrates measurable
traceability coverage; it does not establish that the 88 wiki pages are
semantically complete or that every sentence is entailed by its sources.

\subsection{RQ2: Detection of Seeded Maintenance Faults}

Agent Wiki detected all 120 injected faults and emitted no additional fault
labels, yielding precision and recall of 1.00 in the constructed fixture. The
existence-only baseline detected the 30 unresolved targets but missed the 90
faults that did not contain a broken explicit link, yielding precision 1.00 and
recall 0.25. Table \ref{tab:faultresults} gives class-level results.

\begin{table}[H]
\centering
\small
\begin{tabular}{lrrr}
\toprule
Fault class & Injected & Agent Wiki & Existence-only \\
\midrule
Missing related target & 30 & 30 & 0 \\
Unresolved related target & 30 & 30 & 30 \\
Missing wiki backlink & 30 & 30 & 0 \\
Explicit follow-up & 30 & 30 & 0 \\
\midrule
Total & 120 & 120 & 30 \\
Recall & --- & 1.00 & 0.25 \\
\bottomrule
\end{tabular}
\caption{Fault-injection detection results. No false positives were observed.}
\label{tab:faultresults}
\end{table}

The result answers RQ2 within the defined fault model: lifecycle-aware checks
identify three classes of structural maintenance failure that link existence
alone cannot represent. The perfect score is expected for checks designed
against these explicit invariants and should be interpreted as mutation
adequacy, not as an estimate of defect detection in arbitrary knowledge bases.

\subsection{RQ3: Runtime Cost}

Table \ref{tab:runtime} reports end-to-end timing. The median increased from
56.8 ms for 104 total nodes to 1,302.3 ms for 2,080 nodes. The 95th percentile
at the largest size was 1,344.9 ms. No closure issues were reported in the
clean fixtures.

\begin{table}[H]
\centering
\small
\begin{tabular}{rrrrrr}
\toprule
Raw & Wiki & Nodes & Edges & Median (ms) & P95 (ms) \\
\midrule
100 & 4 & 104 & 200 & 56.8 & 84.2 \\
500 & 20 & 520 & 1,000 & 314.4 & 342.0 \\
1,000 & 40 & 1,040 & 2,000 & 651.3 & 912.3 \\
2,000 & 80 & 2,080 & 4,000 & 1,302.3 & 1,344.9 \\
\bottomrule
\end{tabular}
\caption{End-to-end scan and closure-check runtime over ten repetitions.}
\label{tab:runtime}
\end{table}

These measurements answer RQ3 for the tested implementation and environment.
At the studied documentation scale, complete local checking remained below two
seconds at the 95th percentile. The data suggest approximately proportional
growth over this range, but the study does not fit or claim an asymptotic model.

\section{Discussion}

\subsection{Evidence Closure as a Software Quality Gate}

The evaluation shows why \texttt{processed} should be treated as a verifiable
state rather than a descriptive label. Three quarters of the seeded faults did
not require a broken link: the source could omit its durable target, the target
could omit evidence, or a follow-up could remain open. These failures are
visible only when lifecycle meaning is added to the graph.

The gate is deliberately small. This favors adoption because it can run in a
local terminal, a coding-agent workflow, or continuous integration without a
model call. It also makes failures actionable: each reason code points to a
specific edit. More ambitious semantic checks can be layered on top, but they
need not replace the deterministic foundation.

\subsection{Relationship to Retrieval Systems}

Agent Wiki should not be interpreted as a competitor to RAG. RAG decides what
evidence to retrieve for a query; evidence closure decides whether the corpus
maintains a minimum structural path between source and synthesis. A retrieval
index can be generated from raw notes, wiki pages, or both. Keeping the file
corpus authoritative makes the index disposable and rebuildable.

The present study therefore avoids an answer-quality comparison with a vector
baseline. Such a comparison would require a query set, relevance judgments,
model and embedding controls, and answer-faithfulness evaluation. Those are
important next steps, but including them without an adequately labeled task set
would weaken rather than strengthen the claim made here.

\subsection{Practical Implications}

For independent developers and small teams, the results support three practical
uses. First, source-grounded knowledge can remain in ordinary files without
giving up automated integrity checks. Second, agents can receive repository-
local operating rules and commands instead of relying on hidden product memory.
Third, public software can be separated from private corpora while derived
metadata and evaluation fixtures remain reproducible.

The runtime result also makes closure suitable as a routine pre-commit or
maintenance check at the studied scale. Larger repositories may require indexed
incremental checking. The current shared scanner reads all Markdown files on
each run, prioritizing simplicity over incremental performance.

\section{Threats to Validity}

\paragraph{Construct validity.}
The dependent variables measure structural traceability, not knowledge
correctness, completeness, usefulness, or user trust. Source URLs and hashes do
not establish source authority. Reciprocal links do not establish entailment.
We restrict conclusions accordingly.

\paragraph{Internal validity.}
The four fault classes correspond directly to implemented closure conditions,
so perfect Agent Wiki recall is not surprising. The purpose is to validate that
the executable checks realize the stated invariant. The existence-only
baseline is intentionally minimal and should not be generalized to mature
knowledge-management products. Fault locations are non-overlapping and do not
test interacting mutations.

\paragraph{External validity.}
The observational dataset covers one technical-documentation domain. The
synthetic fixtures use regular fan-out and small Markdown records. Results may
not generalize to collaborative editing, multilingual notes, very large files,
complex YAML, binary-heavy collections, or repositories with millions of
nodes. Runtime values are machine-dependent.

\paragraph{Reliability and reproducibility.}
The public dataset omits Autodesk text and images to respect third-party rights,
which prevents independent semantic evaluation from the archive alone. It does
retain URLs, hashes, structural features, and edges. The experiment script,
generated fixtures, parameters, raw JSON output, and CSV summaries are public,
allowing the reported structural and runtime measurements to be rerun.

\paragraph{Researcher bias.}
The system and evaluation were produced by a single independent developer.
Fault classes and research questions may reflect the system's design. Public
scripts and explicit limitations reduce but do not remove this risk. An
independent replication and a field study with multiple maintainers are needed.

\section{Conclusion}

This study introduced evidence closure as an executable lifecycle invariant for
agent-maintained Markdown knowledge. Agent Wiki separates raw evidence from
synthesized wiki pages and requires processed sources to declare resolved
targets, retain a reciprocal evidence path, and close explicit follow-ups. In a
public graph derived from 1,092 technical-documentation sources, all observable
closure preconditions had complete coverage. In a controlled 120-fault
experiment, evidence closure detected all four fault classes, while an
existence-only checker detected only unresolved targets. End-to-end checking
remained below two seconds at the 95th percentile for a 2,080-node fixture on
the evaluation machine.

The contribution is not a claim that local Markdown replaces retrieval systems
or that structural integrity guarantees truth. It is a smaller systems claim:
durable agent knowledge can expose a deterministic, inexpensive, and
reviewable path from evidence to synthesis. Future work should evaluate
claim-level faithfulness, retrieval quality, maintenance effort, and multi-user
operation on independently labeled corpora.

\section*{Data and Software Availability}

Agent Wiki release \texttt{v0.2.1} and the public FlexSim-derived structural
dataset are available at \url{https://github.com/NimaChu/agent-wiki} and
archived at \url{https://doi.org/10.5281/zenodo.21366895}. The executable IST
evaluation script, machine-readable results, and submission materials are
archived in release \texttt{v0.2.2} at
\url{https://doi.org/10.5281/zenodo.21368436}. The software is licensed under
the MIT License. The original dataset schema, selection, annotations, derived
features, and graph structure are licensed under CC BY 4.0. Autodesk-provided
documentation text, snapshots, and images are not redistributed.

\section*{CRediT Author Statement}

Zijun Chu: Conceptualization, methodology, software, validation, investigation,
data curation, formal analysis, writing--original draft, writing--review and
editing, visualization, and project administration.

\section*{Declaration of Competing Interest}

The author declares no known competing financial interests or personal
relationships that could have appeared to influence the work reported in this
paper.

\section*{Funding}

This research did not receive any specific grant from funding agencies in the
public, commercial, or not-for-profit sectors.

\section*{Declaration of Generative AI and AI-Assisted Technologies in the Manuscript Preparation Process}

During the preparation of this work, the author used OpenAI Codex to assist
with manuscript organization, language drafting, LaTeX preparation, and the
implementation of reproducible evaluation scripts. The author reviewed and
verified the code, data, references, analysis, and manuscript content, edited
the material as needed, and takes full responsibility for the content of the
article.

\section*{Acknowledgements}

The author thanks the open-source communities whose tools and documentation
support local-first software development and reproducible research.

\begin{thebibliography}{99}

\bibitem{yao2023react}
S. Yao, J. Zhao, D. Yu, N. Du, I. Shafran, K. Narasimhan, Y. Cao,
ReAct: Synergizing reasoning and acting in language models,
International Conference on Learning Representations (2023).
\url{https://arxiv.org/abs/2210.03629}.

\bibitem{jimenez2024swebench}
C.E. Jimenez, J. Yang, A. Wettig, S. Yao, K. Pei, O. Press,
K. Narasimhan, SWE-bench: Can language models resolve real-world GitHub
issues?, International Conference on Learning Representations (2024).
\url{https://arxiv.org/abs/2310.06770}.

\bibitem{yang2024sweagent}
J. Yang, C.E. Jimenez, A. Wettig, K. Lieret, S. Yao, K. Narasimhan,
O. Press, SWE-agent: Agent-computer interfaces enable automated software
engineering, Advances in Neural Information Processing Systems 37 (2024).
\url{https://arxiv.org/abs/2405.15793}.

\bibitem{lewis2020rag}
P. Lewis, E. Perez, A. Piktus, F. Petroni, V. Karpukhin, N. Goyal,
H. Kuttler, M. Lewis, W.-t. Yih, T. Rocktaschel, S. Riedel, D. Kiela,
Retrieval-augmented generation for knowledge-intensive NLP tasks,
Advances in Neural Information Processing Systems 33 (2020) 9459--9474.
\url{https://arxiv.org/abs/2005.11401}.

\bibitem{edge2024graphrag}
D. Edge, H. Trinh, N. Cheng, J. Bradley, A. Chao, A. Mody, S. Truitt,
D. Metropolitansky, R.O. Ness, J. Larson, From local to global: A graph RAG
approach to query-focused summarization, arXiv:2404.16130 (2024).
\url{https://arxiv.org/abs/2404.16130}.

\bibitem{packer2023memgpt}
C. Packer, S. Wooders, K. Lin, V. Fang, S.G. Patil, I. Stoica,
J.E. Gonzalez, MemGPT: Towards LLMs as operating systems,
arXiv:2310.08560 (2023).
\url{https://arxiv.org/abs/2310.08560}.

\bibitem{kleppmann2019localfirst}
M. Kleppmann, A. Wiggins, P. van Hardenberg, M. McGranaghan,
Local-first software: You own your data, in spite of the cloud,
in: Proceedings of the 2019 ACM SIGPLAN International Symposium on New Ideas,
New Paradigms, and Reflections on Programming and Software, 2019, pp. 154--178.
\url{https://doi.org/10.1145/3359591.3359737}.

\bibitem{gotel1994trace}
O.C.Z. Gotel, C.W. Finkelstein, An analysis of the requirements traceability
problem, in: Proceedings of the First International Conference on Requirements
Engineering, 1994, pp. 94--101.
\url{https://doi.org/10.1109/ICRE.1994.292398}.

\bibitem{w3cprov}
Y. Gil, S. Miles (Eds.), PROV model primer, W3C Working Group Note,
World Wide Web Consortium, 2013.
\url{https://www.w3.org/TR/prov-primer/}.

\bibitem{hogan2021knowledgegraphs}
A. Hogan, E. Blomqvist, M. Cochez, C. d'Amato, G. de Melo, C. Gutierrez,
S. Kirrane, J.E.L. Gayo, R. Navigli, S. Neumaier, A.-C.N. Ngomo,
A. Polleres, S.M. Rashid, A. Rula, L. Schmelzeisen, J. Sequeda,
S. Staab, A. Zimmermann, Knowledge graphs,
ACM Computing Surveys 54(4) (2021) 1--37.
\url{https://doi.org/10.1145/3447772}.

\bibitem{matuschakEvergreen}
A. Matuschak, Evergreen notes, 2026.
\url{https://notes.andymatuschak.org/} (accessed 15 July 2026).

\bibitem{obsidian}
Obsidian, Obsidian documentation, 2026.
\url{https://help.obsidian.md/} (accessed 15 July 2026).

\bibitem{logseq}
Logseq, Logseq documentation, 2026.
\url{https://docs.logseq.com/} (accessed 15 July 2026).

\end{thebibliography}

\end{document}
